\documentclass[10pt,letterpaper]{article}
\usepackage[T1]{fontenc}
\usepackage[utf8]{inputenc}
\usepackage[spanish,es-tabla]{babel}
\usepackage{csquotes}
\usepackage{mathpazo}
\usepackage[scaled=0.9]{helvet}
\usepackage{courier}
\usepackage{calc}
\usepackage[top=2cm, bottom=2cm, left=2cm, right=2cm]{geometry}
\usepackage{xcolor}
\usepackage{booktabs}
\usepackage{longtable}
\usepackage{array}
\usepackage{ragged2e}
\usepackage{xurl}
\usepackage[strings]{underscore}
\usepackage{fancyhdr}
\usepackage{sectsty}
\usepackage{tocloft}
\usepackage[colorlinks=true,linkcolor=blue,urlcolor=blue,citecolor=blue,breaklinks=true]{hyperref}
\usepackage{enumitem}

\definecolor{uncpblue}{HTML}{1A237E}
\allsectionsfont{\color{uncpblue}}

\renewcommand{\cftsecleader}{\cftdotfill{\cftsecdotsep}}
\cftpagenumbersoff{section}
\cftpagenumbersoff{subsection}
\cftpagenumbersoff{subsubsection}
\renewcommand{\cftsecfont}{\normalfont}
\renewcommand{\cftsubsecfont}{\normalfont}
\renewcommand{\cftsubsubsecfont}{\normalfont}
\setlength{\parindent}{0pt}
\setlength{\parskip}{4pt}
\setlength{\tabcolsep}{3pt}
\renewcommand{\arraystretch}{0.85}

\pagestyle{fancy}
\fancyhf{}
\fancyhead[L]{\small\textcolor{uncpblue}{\textbf{SGSI-UNCP}}}
\fancyhead[R]{\small\textcolor{uncpblue}{D-SGSI-07}}
\fancyfoot[C]{\thepage}
\renewcommand{\headrulewidth}{0.4pt}
\renewcommand{\footrulewidth}{0.4pt}
\setlength{\headheight}{14pt}

\setlist{nosep,leftmargin=*,after=\vspace{2pt}}
\setlength{\emergencystretch}{2em}
\hyphenation{Planea-miento Comuni-caciones}

\newcommand{\makecover}{
\begin{titlepage}
\centering
\vspace*{2cm}
{\Huge\bfseries\color{uncpblue} SGSI-UNCP\par}
\vspace{0.7cm}
{\LARGE Mapa de Procesos del Sistema de Gestión de Seguridad de la
Información\par}
\vspace{0.5cm}
{\large Documento Base del Sistema de Gestión\\de Seguridad de la Información\par}
\vspace{1.5cm}
{\small Universidad Nacional del Centro del Perú\par}
\vfill
\end{titlepage}
}

\begin{document}

\makecover
\setcounter{tocdepth}{2}
\RaggedRight
\tableofcontents
\newpage

\RaggedRight
\textbf{Organización:} Universidad Nacional del Centro del Perú (UNCP)

\textbf{Referencia Normativa:} ISO/IEC 27001:2022 --- Cláusula 4.4
(Sistema de Gestión y sus Procesos), Cláusula 4.1 (Comprensión de la
organización)

\textbf{Documentos Relacionados:} D-SGSI-01 (Contexto Estratégico),
D-SGSI-02 (Alcance del SGSI), D-SGSI-06 (Marco Conceptual), R-SGSI-01
(Inventario de Activos), R-SGSI-02 (Matriz de Riesgos)

\textbf{Mapa de Procesos Institucional:} Resoluciones R. N° 3524-R-2026
(Nivel 0) y R. N° 3526-R-2026 (Nivel 01)

\section{Introducción}\label{introducciuxf3n}

El presente documento describe los procesos institucionales de la UNCP
desde la perspectiva del Sistema de Gestión de Seguridad de la
Información. Cada proceso se analiza identificando los activos de
información críticos que gestiona, los riesgos de seguridad asociados,
los sistemas que lo soportan y los controles aplicables según ISO/IEC
27001:2022.

El mapa de procesos constituye la base para la identificación de activos
de información (R-SGSI-01), la evaluación de riesgos (R-SGSI-02) y la
determinación del alcance del SGSI (D-SGSI-02). La meta del PGTD-UNCP es
incorporar el 90\% de los procesos institucionales al SGSI al año 2028.

\section{Clasificación de Procesos para la
Seguridad}\label{clasificaciuxf3n-de-procesos-para-la-seguridad}

\subsection{Procesos Estratégicos}\label{procesos-estratuxe9gicos}

Los procesos estratégicos definen la dirección, las políticas y la
supervisión del SGSI. Su afectación compromete la gobernabilidad de la
UNCP y la continuidad del sistema.

{\setlength{\RaggedRightRightskip}{0pt plus 5em}\small\sloppy \begin{longtable}[]{@{}
  >{\hspace{0pt}\RaggedRight\arraybackslash}p{(\linewidth - 10\tabcolsep) * \real{0.0826}}
  >{\hspace{0pt}\RaggedRight\arraybackslash}p{(\linewidth - 10\tabcolsep) * \real{0.1193}}
  >{\hspace{0pt}\RaggedRight\arraybackslash}p{(\linewidth - 10\tabcolsep) * \real{0.2661}}
  >{\hspace{0pt}\RaggedRight\arraybackslash}p{(\linewidth - 10\tabcolsep) * \real{0.1651}}
  >{\hspace{0pt}\RaggedRight\arraybackslash}p{(\linewidth - 10\tabcolsep) * \real{0.1927}}
  >{\hspace{0pt}\RaggedRight\arraybackslash}p{(\linewidth - 10\tabcolsep) * \real{0.1743}}@{}}
\toprule\noalign{}
\begin{minipage}[b]{\linewidth}\raggedright
Proceso
\end{minipage} & \begin{minipage}[b]{\linewidth}\raggedright
Descripción
\end{minipage} & \begin{minipage}[b]{\linewidth}\raggedright
Activo de Información Clave
\end{minipage} & \begin{minipage}[b]{\linewidth}\raggedright
Sistema Asociado
\end{minipage} & \begin{minipage}[b]{\linewidth}\raggedright
Riesgo de Seguridad
\end{minipage} & \begin{minipage}[b]{\linewidth}\raggedright
Control ISO 27001
\end{minipage} \\
\midrule\noalign{}
\endhead
\bottomrule\noalign{}
\endlastfoot
Gobierno Digital & Liderazgo del Comité de Gobierno Digital (CGD),
aprobación de políticas, supervisión del PGTD y del SGSI & Actas del
CGD, resoluciones rectorales, políticas de seguridad, informes de avance
& Gestor documental, Microsoft 365 & Alto --- retraso en la toma de
decisiones, pérdida de trazabilidad de acuerdos & A.5.1, A.5.2, A.5.4 \\
Planeamiento Estratégico & Alineamiento del SGSI con el PEI 2024-2030,
PGTD 2026-2030 y presupuesto institucional & PEI, PGTD, informes de
avance de proyectos, matriz de indicadores & Gestor documental,
Microsoft 365 & Alto --- desalineación estratégica, incumplimiento de
metas PGTD & A.5.8, A.5.11 \\
Gestión de Riesgos Institucionales & Aplicación de ISO 31000,
identificación y tratamiento de riesgos institucionales, definición del
apetito de riesgo & Matriz de riesgos institucionales, Declaración de
Aplicabilidad (SoA), planes de tratamiento & R-SGSI-02, D-SGSI-04,
D-SGSI-05 & Alto --- riesgos no mitigados, pérdida de activos críticos &
A.5.6, A.5.7 \\
Gestión de Calidad y Mejora Continua & Auditorías internas del SGSI,
revisión por la dirección, acciones correctivas y preventivas & Informes
de auditoría, actas de revisión por la dirección, RAC, plan de mejora &
P-SGSI-09, R-SGSI-04, F-SGSI-04 & Medio --- pérdida de trazabilidad de
no conformidades, estancamiento del SGSI & A.5.25, A.5.36 \\
Gestión de Cumplimiento Normativo & Seguimiento de obligaciones legales
(Ley N° 29733, D.L. N° 1412, Marco de Confianza Digital), atención a
entes de control (CGR) & Matriz de requisitos legales, informes a la
CGR, registro de sanciones & Gestor documental, Microsoft 365 & Alto ---
sanciones administrativas, multas por incumplimiento de protección de
datos & A.5.31, A.5.32 \\
\end{longtable}\normalsize}

\subsection{Procesos Misionales}\label{procesos-misionales}

Los procesos misionales generan valor directo a la sociedad. Su
interrupción afecta la formación de más de 10,000 estudiantes, la
investigación científica y la responsabilidad social universitaria.

{\setlength{\RaggedRightRightskip}{0pt plus 5em}\small\sloppy \begin{longtable}[]{@{}
  >{\hspace{0pt}\RaggedRight\arraybackslash}p{(\linewidth - 8\tabcolsep) * \real{0.0968}}
  >{\hspace{0pt}\RaggedRight\arraybackslash}p{(\linewidth - 8\tabcolsep) * \real{0.1398}}
  >{\hspace{0pt}\RaggedRight\arraybackslash}p{(\linewidth - 8\tabcolsep) * \real{0.3118}}
  >{\hspace{0pt}\RaggedRight\arraybackslash}p{(\linewidth - 8\tabcolsep) * \real{0.1935}}
  >{\hspace{0pt}\RaggedRight\arraybackslash}p{(\linewidth - 8\tabcolsep) * \real{0.2581}}@{}}
\toprule\noalign{}
\begin{minipage}[b]{\linewidth}\raggedright
Proceso
\end{minipage} & \begin{minipage}[b]{\linewidth}\raggedright
Descripción
\end{minipage} & \begin{minipage}[b]{\linewidth}\raggedright
Activo de Información Clave
\end{minipage} & \begin{minipage}[b]{\linewidth}\raggedright
Sistema Asociado
\end{minipage} & \begin{minipage}[b]{\linewidth}\raggedright
Flujo de Datos Crítico
\end{minipage} \\
\midrule\noalign{}
\endhead
\bottomrule\noalign{}
\endlastfoot
Gestión de Admisión & Proceso de admisión, inscripción de postulantes,
aplicación de examen de admisión, publicación de resultados & Bases de
datos de postulantes, resultados de examen, cuadros de mérito & ERP
ADESA, Portal web, Sistema de admisión & Postulante → inscripción →
examen → calificación → publicación \\
Gestión de Matrícula & Registro de matrícula, asignación de asignaturas,
generación de horarios, registro de estudiantes & Registros de
matrícula, fichas de estudiante, horarios, silabos & ERP ADESA, Portal
del Estudiante & Postulante admitido → registro → asignación de cursos →
horario \\
Gestión de Formación y Notas & Registro de notas, control de asistencia,
gestión de planes de estudio, evaluaciones & Actas de notas, registros
de asistencia, planes de estudio, silabos & ERP ADESA, Moodle 4.1,
Microsoft Teams & Docente → registro de notas → actas → validación →
cierre \\
Gestión de Grados y Títulos & Emisión de grados y títulos, registro de
egresados, trámite ante SUNEDU & Expedientes de grados, resoluciones de
otorgamiento, diplomas, registros SUNEDU & ERP ADESA, Repositorio
DSpace, Portal SUNEDU & Egresado → solicitud → validación académica →
resolución → SUNEDU \\
Gestión de Investigación & Desarrollo de tesis, proyectos de
investigación, patentes, publicaciones científicas & Tesis, papers,
datos de investigación, patentes, registros de proyectos & DSpace,
Sistema de investigación, Microsoft 365 & Investigador → proyecto →
recolección de datos → análisis → publicación → repositorio \\
Gestión de Responsabilidad Social & Proyección social, extensión
universitaria, voluntariado, convenios institucionales & Convenios,
informes de proyección social, registros de beneficiarios & Portal web,
Gestor documental & Comunidad → proyecto → ejecución → informe →
rendición de cuentas \\
Servicios Estudiantiles & Comedor universitario, Centro Médico,
bienestar universitario, actividades culturales y deportivas, becas &
Historias clínicas, registros de becas, fichas socioeconómicas,
registros de beneficiarios & Sistema de bienestar universitario &
Estudiante → solicitud → evaluación → beneficio → seguimiento \\
\end{longtable}\normalsize}

\subsubsection{Flujo de Datos del Proceso Misional Crítico: Gestión
Académica}\label{flujo-de-datos-del-proceso-misional-cruxedtico-gestiuxf3n-acaduxe9mica}

El flujo de datos del proceso misional principal (Gestión Académica) se
compone de las siguientes etapas secuenciales:

{\setlength{\RaggedRightRightskip}{0pt plus 5em}\small\sloppy \begin{longtable}[]{@{}
  >{\hspace{0pt}\RaggedRight\arraybackslash}p{(\linewidth - 8\tabcolsep) * \real{0.0909}}
  >{\hspace{0pt}\RaggedRight\arraybackslash}p{(\linewidth - 8\tabcolsep) * \real{0.1688}}
  >{\hspace{0pt}\RaggedRight\arraybackslash}p{(\linewidth - 8\tabcolsep) * \real{0.2078}}
  >{\hspace{0pt}\RaggedRight\arraybackslash}p{(\linewidth - 8\tabcolsep) * \real{0.1169}}
  >{\hspace{0pt}\RaggedRight\arraybackslash}p{(\linewidth - 8\tabcolsep) * \real{0.4156}}@{}}
\toprule\noalign{}
\begin{minipage}[b]{\linewidth}\raggedright
Etapa
\end{minipage} & \begin{minipage}[b]{\linewidth}\raggedright
Descripción
\end{minipage} & \begin{minipage}[b]{\linewidth}\raggedright
Datos Críticos
\end{minipage} & \begin{minipage}[b]{\linewidth}\raggedright
Sistema
\end{minipage} & \begin{minipage}[b]{\linewidth}\raggedright
Control de Seguridad Aplicable
\end{minipage} \\
\midrule\noalign{}
\endhead
\bottomrule\noalign{}
\endlastfoot
1. Admisión & Inscripción de postulantes, aplicación del examen,
publicación de resultados & Datos personales de postulantes (nombres,
DNI, domicilio), resultados de examen, cuadro de méritos & Portal web,
ERP ADESA & A.9.4 (Restricción de acceso), A.10.1 (Cifrado en tránsito),
A.8.24 (Cifrado en reposo) \\
2. Matrícula & Registro del estudiante admitido, asignación de
asignaturas, generación de horarios & Ficha de estudiante, registro de
matrícula, asignaturas inscritas & ERP ADESA, Portal del Estudiante &
A.9.2 (Gestión de acceso), A.9.3 (Responsabilidades de acceso), A.8.13
(Respaldos) \\
3. Formación y Notas & Registro de notas parciales y finales, control de
asistencia, evaluación docente & Actas de notas, registros de
asistencia, evaluaciones parciales y finales & ERP ADESA, Moodle 4.1 &
A.8.13 (Integridad de respaldos), A.8.2 (Privilegios de acceso), A.12.4
(Registro de eventos) \\
4. Grados y Títulos & Validación de egreso, emisión de grados, registro
ante SUNEDU & Expediente de grado, acta de sustentación, resolución de
otorgamiento, diploma & ERP ADESA, SUNEDU & A.5.32 (Propiedad
intelectual), A.8.2 (Privilegios), A.13.1 (Seguridad de red) \\
5. Archivo y Custodia & Archivamiento de expedientes físicos y
digitales, custodia de actas y resoluciones & Actas originales,
expedientes, resoluciones (físicas y digitales) & Archivo central,
DSpace & A.11.1 (Seguridad física), A.11.2.7 (Eliminación segura), A.8.3
(Manejo de medios) \\
\end{longtable}\normalsize}

\textbf{Nota:} El sistema ERP ADESA (23 módulos) es el habilitador
tecnológico transversal del flujo académico. Su disponibilidad es
crítica: una interrupción prolongada detiene los procesos de matrícula,
registro de notas y emisión de grados, con impacto directo en más de
10,000 estudiantes.

\subsection{Procesos de Apoyo}\label{procesos-de-apoyo}

Los procesos de apoyo son necesarios para que los procesos misionales
funcionen de forma segura y eficiente.

{\setlength{\RaggedRightRightskip}{0pt plus 5em}\small\sloppy \begin{longtable}[]{@{}
  >{\hspace{0pt}\RaggedRight\arraybackslash}p{(\linewidth - 8\tabcolsep) * \real{0.1000}}
  >{\hspace{0pt}\RaggedRight\arraybackslash}p{(\linewidth - 8\tabcolsep) * \real{0.1444}}
  >{\hspace{0pt}\RaggedRight\arraybackslash}p{(\linewidth - 8\tabcolsep) * \real{0.3222}}
  >{\hspace{0pt}\RaggedRight\arraybackslash}p{(\linewidth - 8\tabcolsep) * \real{0.2000}}
  >{\hspace{0pt}\RaggedRight\arraybackslash}p{(\linewidth - 8\tabcolsep) * \real{0.2333}}@{}}
\toprule\noalign{}
\begin{minipage}[b]{\linewidth}\raggedright
Proceso
\end{minipage} & \begin{minipage}[b]{\linewidth}\raggedright
Descripción
\end{minipage} & \begin{minipage}[b]{\linewidth}\raggedright
Activo de Información Clave
\end{minipage} & \begin{minipage}[b]{\linewidth}\raggedright
Sistema Asociado
\end{minipage} & \begin{minipage}[b]{\linewidth}\raggedright
Dependencia Crítica
\end{minipage} \\
\midrule\noalign{}
\endhead
\bottomrule\noalign{}
\endlastfoot
Tecnologías de la Información (OTI) & Gestión de infraestructura TI,
redes, cloud, seguridad, CSIRT, soporte técnico. Solo 2 profesionales
para \textasciitilde11,700 usuarios & Configuraciones de red,
credenciales de administración, logs de seguridad, inventario de activos
& SIEM (por implementar), Microsoft 365 Admin, Huawei Cloud Console &
Crítica --- todos los procesos dependen de TI \\
Gestión de Recursos Humanos & Contratación, capacitación, evaluación de
desempeño, desvinculación, planillas & Expedientes de personal,
planillas, contratos, evaluaciones, declaraciones juradas & SIGA/SIAF
--- Módulo RRHH & Alta --- fuga de datos personales de
\textasciitilde1,700 trabajadores \\
Gestión Financiera y Abastecimiento & Contabilidad, tesorería,
presupuesto, adquisiciones, proveedores & Registros contables, órdenes
de compra, facturas, contratos, PAC & SIGA/SIAF, ERP ADESA (módulos
financieros) & Crítica --- fraude financiero, indisponibilidad de
pagos \\
Gestión Documentaria & Mesa de partes, archivo central, gestión de
expedientes digitales y notificaciones & Documentos recibidos y
emitidos, resoluciones, memorandos, TUPA & Sistema de trámite
documentario (GESDOC) & Alta --- pérdida de trazabilidad documental \\
Asesoría Jurídica & Cumplimiento legal, protección de datos personales,
contratos, defensa institucional & Contratos, resoluciones, directivas,
informes legales, procesos judiciales & Gestor documental & Alta ---
responsabilidad legal por incumplimiento normativo \\
Comunicaciones e Imagen & Gestión de comunicación, portal web, redes
sociales y atención al ciudadano & Contenido web, comunicados oficiales,
cuentas de redes sociales & Portal web, gestor de contenidos & Media ---
suplantación de identidad institucional, desinformación \\
\end{longtable}\normalsize}

\section{Matriz de Dependencias entre
Procesos}\label{matriz-de-dependencias-entre-procesos}

{\setlength{\RaggedRightRightskip}{0pt plus 5em}\small\sloppy \begin{longtable}[]{@{}
  >{\hspace{0pt}\RaggedRight\arraybackslash}p{(\linewidth - 8\tabcolsep) * \real{0.0918}}
  >{\hspace{0pt}\RaggedRight\arraybackslash}p{(\linewidth - 8\tabcolsep) * \real{0.2143}}
  >{\hspace{0pt}\RaggedRight\arraybackslash}p{(\linewidth - 8\tabcolsep) * \real{0.1939}}
  >{\hspace{0pt}\RaggedRight\arraybackslash}p{(\linewidth - 8\tabcolsep) * \real{0.2143}}
  >{\hspace{0pt}\RaggedRight\arraybackslash}p{(\linewidth - 8\tabcolsep) * \real{0.2857}}@{}}
\toprule\noalign{}
\begin{minipage}[b]{\linewidth}\raggedright
Proceso
\end{minipage} & \begin{minipage}[b]{\linewidth}\raggedright
Depende de (Entrada)
\end{minipage} & \begin{minipage}[b]{\linewidth}\raggedright
Provee a (Salida)
\end{minipage} & \begin{minipage}[b]{\linewidth}\raggedright
Sistema Transversal
\end{minipage} & \begin{minipage}[b]{\linewidth}\raggedright
Impacto de Indisponibilidad
\end{minipage} \\
\midrule\noalign{}
\endhead
\bottomrule\noalign{}
\endlastfoot
Admisión & RRHH (convocatoria de personal), Abastecimiento (bienes y
servicios) & Matrícula, SUNEDU (registro de ingresantes) & ERP ADESA,
Portal web & Retraso del ciclo de admisión anual \\
Matrícula & Admisión, Servicios Estudiantiles (becas), Tesorería (pagos)
& Formación (listas de clase), Tesorería (cobros) & ERP ADESA &
Imposibilidad de iniciar el semestre académico \\
Formación y Notas & Matrícula, RRHH (asignación docente) & Grados y
Títulos (egreso), Archivo (actas) & ERP ADESA, Moodle 4.1 & Pérdida de
registros de notas, imposibilidad de emitir grados \\
Grados y Títulos & Formación y Notas, Asesoría Jurídica (legalidad del
egreso) & SUNEDU (registro de egresados), Archivo Central & ERP ADESA,
DSpace & Imposibilidad de emitir diplomas, incumplimiento SUNEDU \\
Investigación & RRHH, Abastecimiento & Repositorio DSpace, Publicaciones
científicas & DSpace, Laboratorios & Pérdida de propiedad intelectual,
retraso en publicaciones \\
Tesorería & Matrícula (pagos), Abastecimiento (proveedores), RRHH
(planillas) & Contabilidad, SUNAT & SIGA/SIAF & Imposibilidad de pagar
planillas y proveedores \\
OTI & Abastecimiento (hardware, licencias, servicios cloud) & Todos los
procesos (conectividad, sistemas, seguridad) & Microsoft 365, Huawei
Cloud, Red de datos & Detención de todos los procesos institucionales \\
\end{longtable}\normalsize}

\textbf{Impacto transversal:} El proceso de OTI es el habilitador
crítico de todos los demás procesos. Con solo 2 profesionales para
\textasciitilde11,700 usuarios, constituye el cuello de botella más
significativo para la implementación del SGSI.

\section{Controles de Seguridad por Nivel de
Proceso}\label{controles-de-seguridad-por-nivel-de-proceso}

{\setlength{\RaggedRightRightskip}{0pt plus 5em}\small\sloppy \begin{longtable}[]{@{}
  >{\hspace{0pt}\RaggedRight\arraybackslash}p{(\linewidth - 4\tabcolsep) * \real{0.2400}}
  >{\hspace{0pt}\RaggedRight\arraybackslash}p{(\linewidth - 4\tabcolsep) * \real{0.4800}}
  >{\hspace{0pt}\RaggedRight\arraybackslash}p{(\linewidth - 4\tabcolsep) * \real{0.2800}}@{}}
\toprule\noalign{}
\begin{minipage}[b]{\linewidth}\raggedright
Nivel de Proceso
\end{minipage} & \begin{minipage}[b]{\linewidth}\raggedright
Controles Prioritarios (ISO 27002)
\end{minipage} & \begin{minipage}[b]{\linewidth}\raggedright
Enfoque de Seguridad
\end{minipage} \\
\midrule\noalign{}
\endhead
\bottomrule\noalign{}
\endlastfoot
\textbf{Estratégico} & A.5.1 (Políticas de seguridad), A.5.2 (Roles y
responsabilidades), A.5.4 (Responsabilidades de la dirección), A.5.8
(Gestión de proyectos), A.5.24 (Gestión de incidentes), A.5.25
(Requisitos legales) & Gobernanza, cumplimiento normativo, supervisión
estratégica, continuidad del programa SGSI \\
\textbf{Misional} & A.8.2 (Privilegios de acceso), A.8.13 (Respaldos),
A.8.20 (Seguridad de redes), A.8.24 (Cifrado), A.5.30 (Continuidad),
A.5.32 (Propiedad intelectual) & Disponibilidad e integridad de datos
críticos, protección de propiedad intelectual, continuidad operativa \\
\textbf{Apoyo} & A.8.8 (Gestión de vulnerabilidades), A.8.16
(Monitoreo), A.8.17 (Servicios de TI), A.5.19 (Proveedores), A.7.1
(Seguridad física), A.5.23 (Seguridad cloud) & Operatividad de
infraestructura, protección de servicios compartidos, gestión de
proveedores \\
\end{longtable}\normalsize}

\section{Responsables de Procesos}\label{responsables-de-procesos}

{\setlength{\RaggedRightRightskip}{0pt plus 5em}\small\sloppy \begin{longtable}[]{@{}
  >{\hspace{0pt}\RaggedRight\arraybackslash}p{(\linewidth - 6\tabcolsep) * \real{0.1731}}
  >{\hspace{0pt}\RaggedRight\arraybackslash}p{(\linewidth - 6\tabcolsep) * \real{0.3654}}
  >{\hspace{0pt}\RaggedRight\arraybackslash}p{(\linewidth - 6\tabcolsep) * \real{0.1538}}
  >{\hspace{0pt}\RaggedRight\arraybackslash}p{(\linewidth - 6\tabcolsep) * \real{0.3077}}@{}}
\toprule\noalign{}
\begin{minipage}[b]{\linewidth}\raggedright
Proceso
\end{minipage} & \begin{minipage}[b]{\linewidth}\raggedright
Dueño del Proceso
\end{minipage} & \begin{minipage}[b]{\linewidth}\raggedright
Unidad
\end{minipage} & \begin{minipage}[b]{\linewidth}\raggedright
Rol en el SGSI
\end{minipage} \\
\midrule\noalign{}
\endhead
\bottomrule\noalign{}
\endlastfoot
Gestión Académica & Vicerrector Académico & Vicerrectorado Académico &
Clasificar y proteger la información académica, autorizar accesos a
sistemas académicos, validar integridad de registros \\
Gestión de Investigación & Vicerrector de Investigación & Vicerrectorado
de Investigación & Proteger la propiedad intelectual, gestionar acceso
al repositorio DSpace, autorizar publicaciones \\
Tecnologías de la Información & Jefe de la OTI & Oficina de Tecnologías
de la Información & Implementar y operar los controles técnicos del
SGSI, administrar infraestructura de seguridad, gestionar incidentes \\
Gestión Financiera & Director de Administración & Dirección de
Administración & Asegurar confidencialidad e integridad de información
financiera, autorizar transacciones críticas \\
Recursos Humanos & Jefe de RRHH & Oficina de Recursos Humanos &
Gestionar confidencialidad de datos personales del personal, controlar
accesos al módulo de planillas \\
Gestión Documentaria & Jefe del Archivo Central & Archivo Central &
Custodiar documentos físicos y digitales, gestionar eliminación segura,
controlar acceso a archivos históricos \\
Asesoría Jurídica & Jefe de Asesoría Jurídica & Oficina de Asesoría
Jurídica & Asegurar cumplimiento legal del SGSI, gestionar protección de
datos personales, revisar contratos de TI \\
\end{longtable}\normalsize}

\section{Mapa de Interacción del
SGSI}\label{mapa-de-interacciuxf3n-del-sgsi}

{\setlength{\RaggedRightRightskip}{0pt plus 5em}\small\sloppy \begin{longtable}[]{@{}
  >{\hspace{0pt}\RaggedRight\arraybackslash}p{(\linewidth - 6\tabcolsep) * \real{0.2267}}
  >{\hspace{0pt}\RaggedRight\arraybackslash}p{(\linewidth - 6\tabcolsep) * \real{0.2933}}
  >{\hspace{0pt}\RaggedRight\arraybackslash}p{(\linewidth - 6\tabcolsep) * \real{0.2133}}
  >{\hspace{0pt}\RaggedRight\arraybackslash}p{(\linewidth - 6\tabcolsep) * \real{0.2667}}@{}}
\toprule\noalign{}
\begin{minipage}[b]{\linewidth}\raggedright
Componente SGSI
\end{minipage} & \begin{minipage}[b]{\linewidth}\raggedright
Procesos Relacionados
\end{minipage} & \begin{minipage}[b]{\linewidth}\raggedright
Documento SGSI
\end{minipage} & \begin{minipage}[b]{\linewidth}\raggedright
Tipo de Interacción
\end{minipage} \\
\midrule\noalign{}
\endhead
\bottomrule\noalign{}
\endlastfoot
Política de Seguridad & Todos los procesos estratégicos & D-SGSI-03 &
Define los principios y responsabilidades de seguridad aplicables a
todos los procesos \\
Evaluación de Riesgos & Todos los procesos misionales y de apoyo &
D-SGSI-04, R-SGSI-02 & Identifica riesgos específicos por proceso y
define tratamientos \\
Control de Acceso & Gestión Académica, Investigación, RRHH, Tesorería &
P-SGSI-03 & Regula quién puede acceder a qué sistemas y datos según su
rol \\
Gestión de Incidentes & OTI, todos los procesos (reporte) & P-SGSI-02 &
Canaliza la detección, reporte y respuesta a incidentes de seguridad \\
Continuidad & Gestión Académica, Investigación, Tesorería & P-SGSI-05 &
Asegura la continuidad de los procesos críticos ante interrupciones \\
Gestión de Cambios & Todos los procesos (cambios en sistemas, personal,
normativa) & P-SGSI-04 & Controla los cambios que puedan afectar la
seguridad de la información \\
Capacitación & RRHH, todos los procesos (concientización) & P-SGSI-08 &
Desarrolla competencias de seguridad en todos los niveles \\
Auditoría Interna & Todos los procesos (evaluación periódica) &
P-SGSI-09 & Verifica la eficacia de los controles en cada proceso \\
Eliminación Segura & Archivo Central, OTI, todos los procesos &
P-SGSI-07 & Gestiona la destrucción segura de información cuando pierde
vigencia \\
\end{longtable}\normalsize}

\section{Documentos Relacionados}\label{documentos-relacionados}

{\setlength{\RaggedRightRightskip}{0pt plus 5em}\small\sloppy \begin{longtable}[]{@{}
  >{\hspace{0pt}\RaggedRight\arraybackslash}p{(\linewidth - 4\tabcolsep) * \real{0.1600}}
  >{\hspace{0pt}\RaggedRight\arraybackslash}p{(\linewidth - 4\tabcolsep) * \real{0.1600}}
  >{\hspace{0pt}\RaggedRight\arraybackslash}p{(\linewidth - 4\tabcolsep) * \real{0.6800}}@{}}
\toprule\noalign{}
\begin{minipage}[b]{\linewidth}\raggedright
Código
\end{minipage} & \begin{minipage}[b]{\linewidth}\raggedright
Nombre
\end{minipage} & \begin{minipage}[b]{\linewidth}\raggedright
Relación con el Mapa de Procesos
\end{minipage} \\
\midrule\noalign{}
\endhead
\bottomrule\noalign{}
\endlastfoot
D-SGSI-01 & Análisis del Contexto Estratégico & Identifica factores
externos e internos que afectan los procesos \\
D-SGSI-02 & Alcance del SGSI & Define qué procesos están dentro del
alcance del SGSI \\
D-SGSI-06 & Marco Conceptual del SGSI & Establece el modelo conceptual y
ciclo PHVA que rige los procesos \\
R-SGSI-01 & Inventario de Activos de Información & Cataloga los activos
de información identificados por proceso \\
R-SGSI-02 & Matriz de Evaluación y Tratamiento de Riesgos & Documenta
los riesgos asociados a cada proceso y sus controles \\
P-SGSI-00 & Control de Documentos y Registros & Regula la documentación
de los procesos del SGSI \\
\end{longtable}\normalsize}

\end{document}
